%!TEX root = ../../main.tex
\chapter{C 分析}

\section{分析类图（边界 / 控制 / 实体三分法）}

分析类图把系统涉及的类按职责分成三类：
\textbf{边界类}是系统与外部系统交互的界面，一个有接口的外部系统映射一个边界类；
只负责收发消息、不做业务判断；\textbf{控制类}负责协调与流程控制，把执行者请求分配给实体类完成，
不持有业务数据，因此不写 get/set 一类的存取操作；\textbf{实体类}是核心域概念，
属性只写直接描述该类特征的领域字段，不画标识属性（主键）与状态属性（状态交由状态机负责）。

\subsection{边界类与控制类}

边界类与控制类分两张图给出：柜台电脑端一侧见图~\ref{fig:05a-boundary-control-class}，微信小程序一侧见图~\ref{fig:05a2-boundary-control-class}——一个外部系统一张图。边界类按\textbf{外部系统}划分而不是按角色划分：
柜台电脑端（老板与销售共用一套系统）与微信小程序（客户、供应商、加工师傅、合作商户共用一个小程序）
各对应一个边界类，合计 2 个；同一外部系统内各角色能看到的字段与按钮由登录角色决定，不另立边界类。
登录在两个外部系统都有入口，因此登录控制被两个边界类共同依赖。
控制类与阶段 B 的用例一一对应，命名一律“用例名 + 控制”，每个控制类只承担“接收请求、组织协作、返回结果”的责任。
边界类到控制类是转发关系，控制类到实体类是分配责任，前者见表~\ref{tab:b2c}，后者见各用例的分析序列图。
两者的方向都是\textbf{依赖}：控制类不保存数据，也不调用实体类的取值/赋值操作。

\figp{分析类图（一）：边界类与柜台端控制类}{05a-boundary-control-class}

\fig{分析类图（二）：微信小程序边界类与控制类}{05a2-boundary-control-class}

\needspace{8\baselineskip}
\begin{table}[H]
\centering
\small
\caption{边界类到控制类的对应关系（依赖）}
\label{tab:b2c}
\begin{tabularx}{\textwidth}{@{}l Y@{}}
\toprule
边界类 & 覆盖的控制类（对应用例） \\
\midrule
柜台端界面 & 登记成交、登记货品入库、登记退货、维护客户信息、登记送检与证书、查看销售与周转统计、维护基础数据、开通外部账户、审核出让申请、登记货款记录、维护调货记录、登记寄售合同、结算寄售货款、结算调货台账、登记分期计划、记录分期还款、查看库存流水、跟进预定到店（共 18 个） \\
微信小程序界面 & 浏览货品、预约查看货品、查看调货价格、预定商品、提交货品出让申请、预约上门时间、查看加工任务、提交完工交付申请、查询预约（共 9 个） \\
两端共用 & 登录（UC-08 被各端用例以 include 引用，在两个外部系统都有入口） \\
\bottomrule
\end{tabularx}
\end{table}


\subsection{实体类及其关联（关联与组合）}

实体类及其关联如图~\ref{fig:05b1-entity-goods}。共 26 个实体类：货品、入库批次、证书、成交、退货、
送检、预定、预约、出让申请、上门预约、加工任务、完工交付申请、货款、调货、寄售合同、寄售明细、
台账结算、分期计划、库存流水、客户、来源方、鉴定机构、货品类别、用户、角色、权限。
来源方是供应商、加工师傅与同行商户的合一：三类的属性与操作完全一致，只以类别取值区分
（与货款按业务类型取值区分同理），类别在同一取值内名称唯一。
实体类的操作列与下节分析序列图的消息一一对应——某条消息传进来，就是该实体类的操作。
类名不带“记录/信息/数据/情况”等后缀，不画标识属性（主键），
也不画名为“状态”的属性（货品状态交由状态机负责）。合同编号、证书编号等业务编码是印在纸质单据上、用户可引用的领域字段，不是主键标识，故保留在类图中；主键（代理键）不画，数据层与业务编码分离。

客户类有一条自反关联：来源渠道含“朋友介绍”，因而客户之间存在介绍关系。
\textbf{普通关联}用于可共享、或同一对对象可多次发生的关系（如同一件货品的多次成交）；\textbf{组合}用于部分不可共享、且不能脱离整体存在的关系，共四处——证书挂货品、完工交付申请挂加工任务、货品入库批次、寄售明细挂寄售合同；控制类到实体类的\textbf{依赖}见分析类图（一）。\textbf{泛化}不采用：各类账户的差别只在权限与登录入口，核心属性完全一致，属于属性值差别；用户与角色为多对多、角色与权限也为多对多，角色与权限各自成类成表，不为五类角色各建子类。
类之间的 33 条关联连同多重性与角色名见表~\ref{tab:rels}。

\figp{实体类之间的关联（一）：货品与入库批次、证书}{05b1-entity-goods}
\figp{实体类之间的关联（二）：货品与成交}{05b2-entity-deal}
\figtwoW{0.72}{实体类之间的关联（三）：送检与证书}{05b3-entity-inspect}{1.0}{实体类之间的关联（四）：预约、货品、预定与客户}{05b4-entity-book}
\figtwo[0.60]{实体类之间的关联（五）：来源方、出让申请与上门预约}{05b5-entity-transfer}{实体类之间的关联（六）：货品、加工任务与来源方}{05b6-entity-process}
\figp{实体类之间的关联（七）：货品类别与货品的分类关系}{05b8-entity-category}
\figp{实体类之间的关联（八）：货品送检与送检单}{05b9-entity-goods-inspect}
\figp{实体类之间的关联（九）：调货与来源方往来}{05b11-entity-money-partner}
\figtwo[0.72]{实体类之间的关联（十）：加工任务与完工交付申请}{05b12-entity-delivery}{实体类之间的关联（十一）：送检、鉴定机构与证书}{05b13-entity-cert-org}
\figp{实体类之间的关联（十二）：货品的调拨}{05b14-entity-goods-transfer2}
\figtwo[0.72]{实体类之间的关联（十三）：成交与客户下单}{05b15-entity-deal-customer}{实体类之间的关联（十四）：成交与退货}{05b16-entity-deal-return}
\figp{实体类之间的关联（十五）：来源方、寄售合同、寄售明细与货品}{05b17-entity-consign}
\figp{实体类之间的关联（十六）：来源方、台账结算、调货与寄售明细}{05b18-entity-settle}
\figp{实体类之间的关联（十七）：成交与分期计划}{05b19-entity-installment}
\figp{实体类之间的关联（十八）：货品、用户与库存流水}{05b20-entity-stocklog}
\figp{实体类之间的关联（十九）：用户、角色与权限}{05b21-entity-rbac}

\needspace{8\baselineskip}
\begin{table}[H]
\centering
\small
\caption{实体类之间的关联（多重性与角色名）}
\label{tab:rels}
\begin{tabularx}{\textwidth}{@{}l c l c Y@{}}
\toprule
源类 & 源端 & 目标类 & 目标端 & 角色名与说明 \\
\midrule
入库批次 & 1 & 货品 & 0..* & 含货品（一件货品只属一个批次）\textbf{（组合）} \\
货品 & 1 & 证书 & 0..* & 挂证书（证书不可共享，货品注销则证书无意义）\textbf{（组合）} \\
加工任务 & 1 & 完工交付申请 & 0..* & 产出（完工交付申请不可脱离任务）\textbf{（组合）} \\
来源方 & 1 & 寄售合同 & 0..* & 寄出（类别=供应商） \\
寄售合同 & 1 & 寄售明细 & 0..* & 含（一份合同多件寄售货品）\textbf{（组合）} \\
货品 & 1 & 寄售明细 & 0..* & 被寄售（一货品同一时间至多一条未结算明细，结算后可再寄售） \\
来源方 & 1 & 台账结算 & 0..* & 结算往来 \\
调货 & 1 & 台账结算 & 0..1 & 结算目标（一笔调货至多一次结算） \\
寄售明细 & 1 & 台账结算 & 0..1 & 结算目标（一件寄售货品至多一次结算） \\
成交 & 1 & 分期计划 & 0..1 & 分期（一份成交至多一份计划） \\
货品 & 1 & 库存流水 & 0..* & 记录（每次状态变更必写流水） \\
用户 & 1 & 库存流水 & 0..* & 经手 \\
用户 & 0..* & 角色 & 0..* & 具有（多对多，落表 t\_user\_role） \\
角色 & 0..* & 权限 & 0..* & 授予（多对多，落表 t\_role\_permission） \\
货品类别 & 1 & 货品 & 0..* & 分类 \\
货品 & 1 & 成交 & 0..* & 售出（退货后仍可再售，故可多次） \\
客户 & 0..1 & 成交 & 0..* & 下单（散客不关联客户档案，允许为空） \\
成交 & 1 & 退货 & 0..* & 被退（退货是后续事件而非组成部分，故不画组合） \\
货品 & 1 & 送检 & 0..* & 送检 \\
送检 & 0..1 & 证书 & 0..* & 出具（证书由送检出具，随货品带入的证书无对应送检） \\
鉴定机构 & 1 & 送检 & 0..* & 送检机构 \\
鉴定机构 & 1 & 证书 & 0..* & 签发机构 \\
货品 & 1 & 预定 & 0..* & 被预定 \\
预定 & 0..* & 客户 & 0..1 & 登记给（预定以姓名与手机号提交，已建档客户才关联） \\
客户 & 0..1 & 客户 & 0..* & 介绍（自反关联：来源渠道含“朋友介绍”） \\
货品 & 1 & 预约 & 0..* & 被预约看 \\
来源方 & 1 & 出让申请 & 0..* & 提交（类别=供应商，含回流出让方） \\
出让申请 & 0..1 & 上门预约 & 0..* & 预约上门 \\
完工交付申请 & 0..1 & 上门预约 & 0..* & 预约上门（UC-15 规则①） \\
货品 & 1 & 加工任务 & 0..* & 加工 \\
来源方 & 1 & 加工任务 & 0..* & 承接（类别=加工师傅） \\
货品 & 1 & 调货 & 0..* & 调拨 \\
来源方 & 1 & 调货 & 0..* & 往来方（类别=同行商户） \\
\bottomrule
\end{tabularx}
\end{table}

\section{分析序列图}

序列图对应每个用例的路径，每个用例的基本路径各画一张图：
简单的扩展点合并到基本路径图，复杂扩展点单独一张并与基本路径图链接。
研究对象是系统内各类，软件类——执行者只发消息到边界类，
边界类把请求转给控制类，控制类再组织实体类完成业务。
消息代表责任分配而非数据流动，控制类只分配责任，不取存实体类的属性。
28 个用例的基本路径逐一给出，如图~\ref{fig:06a-uc01-sale} 至图~\ref{fig:06aa-uc28-booking-follow}。

\figp[1.0]{登记成交（UC-01）分析序列图}{06a-uc01-sale}
\figp[1.0]{登记货品入库（UC-02）分析序列图}{06b-uc02-goods-in}
\figp[1.0]{登记退货（UC-03）分析序列图}{06c-uc03-return}
\figp[1.0]{维护客户信息（UC-04）分析序列图}{06i-uc04-customer}
\figp[1.0]{登记送检与证书（UC-05）分析序列图}{06d-uc05-inspect}
\figp[1.0]{查看销售与周转统计（UC-06）分析序列图}{06j-uc06-stats}
\figp[1.0]{维护基础数据（UC-07）分析序列图}{06k-uc07-basicdata}
\figp[1.0]{登录（UC-08）分析序列图}{06l-uc08-login}
\figp[1.0]{浏览货品（UC-09）分析序列图}{06m-uc09-browse}
\figp[1.0]{预约查看货品（UC-10）分析序列图}{06n-uc10-reserve}
\figp[1.0]{查看调货价格（UC-11）分析序列图}{06o-uc11-transfer-price}
\figp[1.0]{开通外部账户（UC-12）分析序列图}{06p-uc12-open-account}
\figp[1.0]{预定商品（UC-13）分析序列图}{06e-uc13-book}
\figp[1.0]{提交货品出让申请（UC-14）分析序列图}{06f-uc14-transfer-apply}
\figp[1.0]{预约上门时间（UC-15）分析序列图}{06q-uc15-visit-appt}
\figp[1.0]{查看加工任务（UC-16）分析序列图}{06r-uc16-process-task}
\figp[1.0]{提交完工交付申请（UC-17）分析序列图}{06s-uc17-delivery}
\figp[1.0]{审核出让申请（UC-18）分析序列图}{06g-uc18-transfer-review}
\figp[1.0]{登记货款记录（UC-19）分析序列图}{06t-uc19-payment}
\figp[1.0]{维护调货记录（UC-20）分析序列图}{06h-uc20-transfer-rec}
\figp[1.0]{登记寄售合同（UC-21）分析序列图}{06t-uc21-consign}
\figp[1.0]{结算寄售货款（UC-22）分析序列图}{06u-uc22-consign-settle}
\figp[1.0]{结算调货台账（UC-23）分析序列图}{06v-uc23-transfer-settle}
\figp[1.0]{登记分期计划（UC-24）分析序列图}{06w-uc24-installment}
\figp[1.0]{记录分期还款（UC-25）分析序列图}{06x-uc25-installment-pay}
\figp[1.0]{查看库存流水（UC-26）分析序列图}{06y-uc26-stock-log}
\figp[1.0]{查询预约（UC-27）分析序列图}{06z-uc27-resv-query}
\fig[1.0]{跟进预定到店（UC-28）分析序列图}{06aa-uc28-booking-follow}

\section{分析状态机图（只画核心类）}

状态机图只画核心类，每条迁移的触发事件与规约来源一并标注。
共画九张：货品的八状态总览及分主题三张（售出与预留、调货与代销、送检与处置）、出让申请（三状态）、加工任务（四状态）、预定记录（五状态）、寄售合同（三状态）、分期计划（三状态）。
选取依据是“核心类、状态多的事物”：货品是全系统流转最频繁的对象，加工任务是组合结构的根，
预定记录的状态直接驱动额度规则；寄售合同与分期计划是新扩展业务的生命周期载体。

\subsection{货品生命周期（八状态，分三张给出）}

货品是全系统流转最频繁的对象，八状态取自 UC-01 的状态枚举：在柜、已订、已售、预留、已调出、代销中、已送检、待处理。
状态之间的迁移共 18 条，总体结构见图~\ref{fig:07a-goods-state}；
再按业务主题分三张细化：售出与预留（图~\ref{fig:07a1-goods-sale}）、
调货与代销（图~\ref{fig:07a2-goods-transfer}）、送检与处置（图~\ref{fig:07a3-goods-inspect}）。
各迁移及其规约来源见表~\ref{tab:goods-transitions}。

\figp{货品生命周期状态机（八状态总览）}{07a-goods-state}

\figp{货品状态机（一）：售出与预留}{07a1-goods-sale}
\figp{货品状态机（二）：调货与代销}{07a2-goods-transfer}
\figp{货品状态机（三）：送检与处置}{07a3-goods-inspect}

\needspace{8\baselineskip}
\begin{table}[H]
\centering
\small
\caption{货品状态迁移及其规约来源}
\label{tab:goods-transitions}
\begin{tabularx}{\textwidth}{@{}l l Y@{}}
\toprule
迁移 & 触发事件 & 规约来源 \\
\midrule
$[*]\rightarrow$在柜 & 入库建档（采购/回收/调货/加工完工/寄售） & UC-02 \\
在柜 $\rightarrow$ 预留 & 客户预留，或预定确定到店 & UC-01 规则③、UC-13 规则⑥ \\
预留 $\rightarrow$ 在柜 & 取消预留（到店看过暂不决定） & UC-13 规则⑥ \\
在柜 $\rightarrow$ 已售 & 全款成交 & UC-01 \\
在柜 $\rightarrow$ 已订 & 收定金、分期留货 & UC-01 规则④ \\
已订 $\rightarrow$ 已售 & 尾款付清并交付出库 & UC-01 规则④ \\
在柜 $\rightarrow$ 已调出 & 调出给合作商户 & UC-20 \\
已调出 $\rightarrow$ 已售 & 鉴赏期内客户确定留货 & UC-20 规则③ \\
已调出 $\rightarrow$ 在柜 & 退货退回 & UC-20 \\
在柜 $\rightarrow$ 代销中 & 寄卖在同行处 & UC-01 规则⑤ \\
代销中 $\rightarrow$ 已售 & 同行售出、与同行结算 & UC-01 规则⑤ \\
在柜 $\rightarrow$ 已送检 & 送检 & UC-05 \\
预留 $\rightarrow$ 已送检 & 送检 & UC-05 \\
已送检 $\rightarrow$ 在柜 & 证书录入、恢复送检前 & UC-05 \\
已送检 $\rightarrow$ 预留 & 证书录入、恢复送检前 & UC-05 \\
已送检 $\rightarrow$ 待处理 & 检测不合格 & UC-05 扩展 4a \\
已售 $\rightarrow$ 在柜 & 退货（成交后 7 天内、真假/质量） & UC-03 \\
待处理 $\rightarrow$ $[*]$ & 退回供应商或停售 & UC-05 扩展 4a \\
\bottomrule
\end{tabularx}
\end{table}

状态机还体现两条不变式：一是预留、已订期间货品不得再售他人；
二是已调出与代销中的货权不同——已调出货权转移给商户，代销中货权仍归本店。

\subsection{出让申请（三状态）}

出让申请的生命周期如图~\ref{fig:07b-apply-state}：供应商提交后进入“待审核”，
老板审核通过则转“已通过”并转入入库流程，驳回则转“已驳回”。
迁移全部来自 UC-14（提交）与 UC-18（审核）。

\fig{出让申请状态机（三状态）}{07b-apply-state}

\subsection{加工任务（四状态）}

加工任务的状态机如图~\ref{fig:07c-process-task-state}。它是完工交付申请的组合整体，
“在加工”是提交完工交付申请的前置条件，状态迁移直接由 UC-16、UC-17 驱动。

\fig{加工任务状态机（四状态）}{07c-process-task-state}

\subsection{预定记录（五状态）}

预定记录的状态机如图~\ref{fig:07d-booking-state}。额度规则、爽约判定都由状态驱动；
销售在当日到店名单里联系客户后置为“已联系”，客户到店置“已达成”，
超时未到置“已爽约”，两条结束路径都按 UC-13 规则⑤处理额度。
预定本身不锁货，客户确定到店时置为“预留”的是货品状态（见 07a）。

\fig{预定记录状态机（五状态）}{07d-booking-state}

\subsection{寄售合同（三状态）}

寄售合同的状态机如图~\ref{fig:07e-consign-state}。登记寄售合同后即“生效”，
提前终止转“已终止”，到期日过仍未终止转“已到期”；
两种结束态下未结算明细对应的货品都退回供应商。

\fig{寄售合同状态机（三状态）}{07e-consign-state}

\subsection{分期计划（三状态）}

分期计划的状态机如图~\ref{fig:07f-installment-state}。登记分期计划后为“未还”，
到期日过仍有未还转“逾期”，还清全部转“已结清”；
逾期不是终点，补清欠款后同样进入“已结清”。

\figp{分期计划状态机（三状态）}{07f-installment-state}

% ============================================================
% D 阶段：设计
% ============================================================

% ============================================================
% D 阶段：设计
% ============================================================
